\subsection{Formal checking in Lean}\label{sec:collab-lean}

\subsubsection*{Lean and what a formal check establishes}

Lean~4 \citep{demoura2021} is an interactive proof assistant: a programming language in which definitions,
statements and proofs are written in a formal language based on dependent type theory, and in which every proof is
checked by a small trusted program, the kernel. A proof may be written step by step or with the help of tactics,
procedures that search for parts of a proof automatically, but whatever a tactic produces is checked again by the
kernel. Mathlib \citep{mathlib2020} is the community library of formalized mathematics for Lean; it supplies the
real numbers, groups, limits and the other notions used below.

Two features of Lean are used to state precisely what has been checked. A proof can be left unfinished by writing
\texttt{sorry}, which Lean accepts with a warning; a finished development contains none. The command
\texttt{\#print axioms} lists the foundational axioms on which a theorem depends; ordinary classical mathematics in
Lean uses exactly three of them, \texttt{propext} (propositional extensionality), \texttt{Classical.choice} (the
axiom of choice) and \texttt{Quot.sound} (soundness of quotients), and an unfinished proof shows up in this list as
\texttt{sorryAx}.

A successful check establishes that the formal statements follow from the formal definitions. It does not
establish that the formal definitions and statements say what the paper says, which must be checked by a human
reader, and it says nothing about results that enter as hypotheses or are not formalized.

For this paper a formal check is useful for three reasons. First, the proofs were drafted with AI assistance, and
the errors recorded in Table~\ref{tab:errors} show that such drafts can contain steps that read well and are wrong;
the kernel does not depend on how a proof was produced. Second, the steps that are specific to loss ratios
(Lemmas~\ref{lem:premium}--\ref{lem:LI}) are elementary but delicate, with case distinctions at $\Pi(X)=0$, at
break-even books and at the price point, which is where informal checking is most likely to miss a case. Third,
formalization forces the hypotheses to be stated exactly, and here it shows that those steps use only a few
algebraic properties of comonotonicity and equality in law (see ``What was formalized'' below).

\subsubsection*{What was formalized}\label{sec:lean-scope}

The formalization works in an abstract setting. Positions form an additive commutative group $\Xp$ with an additive
embedding $c\mapsto\mathrm{const}(c)$ of the sure losses, a relation ``comonotone'' under which every position is
comonotone with every constant and which is preserved when both positions are shifted by constants, and a relation
``equal in law'' preserved by a common cash shift. A loss ratio is a map from $\Xp$ to $\R\cup\{\mathtt{none}\}$,
where \texttt{none} stands for a non-finite value, so that the labels $\infty$ and $\star$ are identified as in
\axLI. The axioms read as follows; \texttt{prem} is the implied premium and \texttt{rhoD} is $\Pi\cdot\LR$.

{\small
\begin{verbatim}
def AxP : Prop := ∀ x : X, ∃! c : ℝ, LR (x + S.const c) = none
def AxN : Prop := ∀ c : ℝ, c ≠ 0 → LR (S.const c) = some 1
def AxLI : Prop := ∀ x y : X, S.eqd x y → LR x = LR y
def AxCM (hP : AxP S LR) : Prop :=
  ∀ x y : X, S.comon x y → LR x ≠ none → LR y ≠ none →
    prem hP x + prem hP y ≠ 0 →
    LR (x + y) = some ((rhoD hP x + rhoD hP y) / (prem hP x + prem hP y))
\end{verbatim}
}

These coincide with Axioms~\ref{ax:P}, \ref{ax:N}, \ref{ax:CM} and~\ref{ax:LI}. The file has about 500 lines with
10 definitions and 26 lemmas and theorems. Table~\ref{tab:lean} lists the results of the paper that are formalized.

\begin{table}[ht]
\centering\small
\begin{tabular}{@{}ll@{}}
\toprule
\textbf{Result in the paper} & \textbf{Lean name}\\
\midrule
Implied premium, uniqueness and cash additivity (Lemma~\ref{lem:premium}) & \texttt{prem}, \texttt{prem\_shift}, \texttt{prem\_const}\\
Implied risk measure and cash additivity (Lemma~\ref{lem:implied}) & \texttt{rho}, \texttt{rho\_cash}\\
Möbius form and residue (Proposition~\ref{prop:readings}(a)--(c)) & \texttt{mobius}, \texttt{rho\_eq}, \texttt{residue}\\
Comonotonic additivity, also on the break-even set (Lemma~\ref{lem:comon}) & \texttt{comon\_additive}\\
Law invariance of the implied pair (Lemma~\ref{lem:LI}) & \texttt{LI\_transfer}\\
Converse construction and uniqueness of the pair & \texttt{converse}, \texttt{unique\_pair}\\
Association (Corollary~\ref{cor:assoc}(a)) & \texttt{association}\\
Pooling at non-finite values (Corollary~\ref{cor:projpool}) & \texttt{pooling\_extended}\\
Algebraic core of Proposition~\ref{prop:F2} & \texttt{no\_monotone\_core}\\
Theorem~\ref{thm:A}, given the representation step & \texttt{main\_conditional}\\
\bottomrule
\end{tabular}
\caption{Results formalized in Lean.}\label{tab:lean}
\end{table}

The theorem \texttt{main\_conditional} takes the representation step as a hypothesis: any cash-additive,
comonotonic additive, law-invariant and regular functional is of the required form \citep[Thm.~1]{www2020}. From
\axP, \axN, \axCM{} and \axLI{} it derives that the implied premium and the implied risk measure satisfy the
hypotheses of that step and that $\LR=\bracket{\rho}{\Pi}$ everywhere. Not formalized are the representation step
itself (Lemmas~\ref{lem:cashlip}--\ref{lem:rep}), which would require a theory of Choquet integrals that Mathlib does not
provide, the fact that $L^\infty$ with comonotonicity and equality in law is an instance of
the abstract setting, the layer formula, and the examples and counterexamples.

\subsubsection*{How the check was carried out}

The Lean file was written by the main assistant. Lean could not be installed in the assistant's cloud workspace,
whose network policy blocks the servers that distribute Lean, so the file was compiled on the author's computer.
The author installed Lean and created the project; the assistant read the build output from the project folder
through the link between the Claude desktop application and the author's computer, wrote corrected versions of the
file into the same folder, and the author ran the next build (Figure~\ref{fig:leanloop}).

\begin{figure}[ht]
\centering
\begin{tikzpicture}[node distance=6mm and 14mm,
  box/.style={draw,rounded corners,align=center,font=\small,minimum height=10mm,text width=34mm},
  >={Stealth[length=2mm]}]
\node[box] (a) {Assistant writes or\\corrects \texttt{Basic.lean}};
\node[box,right=of a] (b) {File written into the\\author's project folder};
\node[box,below=of b] (c) {Author runs\\\texttt{lake build}};
\node[box,left=of c] (d) {Assistant reads\\\texttt{build\_log.txt}};
\draw[->] (a) -- (b);
\draw[->] (b) -- (c);
\draw[->] (c) -- (d);
\draw[->] (d) -- node[left,font=\footnotesize,align=center]{errors}(a);
\node[box,right=of c,text width=30mm] (e) {Build succeeds;\\axiom report};
\draw[->] (c) -- (e);
\end{tikzpicture}
\caption{The build-and-repair loop used for the Lean check.}\label{fig:leanloop}
\end{figure}

Table~\ref{tab:leantime} gives the timeline. Times are in UTC and are taken from the session records and from the
modification times of the files in the project folder.

\begin{table}[ht]
\centering\small
\begin{tabular}{@{}>{\raggedright\arraybackslash}p{3.3cm}>{\raggedright\arraybackslash}p{11.0cm}@{}}
\toprule
\textbf{Time (UTC)} & \textbf{Step}\\
\midrule
21 Oct, 21:38 & The assistant writes a first Lean file for the results in Table~\ref{tab:lean}; it cannot be compiled
  in the cloud workspace.\\
24 Oct, 17:08 & The author asks how a formal proof check can be used; the assistant describes what is feasible
  (the loss-ratio steps), what is hard (the representation step) and how to install Lean.\\
24 Oct, 17:26--17:31 & The author installs the Lean extension for the code editor and the Lean toolchain manager.\\
24 Oct, 17:35--17:37 & The author creates the project \texttt{loss\_ratio\_lean} with Mathlib; Lean and Mathlib
  version~4.36.0-rc1; the precompiled Mathlib files are downloaded.\\
24 Oct, 17:38--17:39 & The assistant inspects the project configuration, confirms that it is complete, and places
  the Lean file in it, replacing the template placeholder.\\
25 Oct, 12:47--12:53 & First attempt to build: the command \texttt{lake} is not found, because the toolchain had not
  been added to the terminal's search path; fixed by loading the toolchain environment.\\
25 Oct, 12:55 & Build 1: four errors, none in a definition or statement. One proof stopped one algebraic step short;
  in three proofs Lean's simplifier rewrote $\mathrm{const}(-c)$ as $-\mathrm{const}(c)$, so that cash additivity
  no longer applied.\\
25 Oct, 12:57 & The assistant writes a corrected file and appends \texttt{\#print axioms} for ten main theorems.\\
25 Oct, 13:04 & Build 2: one error, in the converse construction, where the goal was displayed in an unexpanded
  form. The axiom report shows the three standard axioms for nine theorems and \texttt{sorryAx} for the converse.\\
25 Oct, 13:05 & The assistant adds one line that states the goal explicitly.\\
25 Oct, 13:08 & Build 3 succeeds. All ten theorems depend only on the three standard axioms.\\
25 Oct, 13:09 & The assistant compares the Lean definitions of \axP, \axN, \axCM{} and \axLI{} and the theorem
  statements with Section~\ref{sec:axioms}.\\
25 Oct, 13:37 & At the author's request, the project's own build output is deleted and the file is rebuilt from
  scratch: success in 8.8\,s, with the same axiom report.\\
\bottomrule
\end{tabular}
\caption{Timeline of the Lean check.}\label{tab:leantime}
\end{table}

\subsubsection*{Checks performed}

The following checks were made to ensure that the outcome is what is reported.

\begin{itemize}[nosep]
\item \emph{No change to statements during repair.} The corrections after builds 1 and 2 changed proof steps only;
  no definition and no theorem statement was altered.
\item \emph{No unfinished proofs.} The final file contains no \texttt{sorry}, no \texttt{admit} and no added
  \texttt{axiom}.
\item \emph{Axiom report.} \texttt{\#print axioms} for each of the ten theorems in Table~\ref{tab:lean} lists only
  \texttt{propext}, \texttt{Classical.choice} and \texttt{Quot.sound}.
\item \emph{Clean rebuild.} The final build started from an empty build directory, so its result does not rely on
  cached output.
\item \emph{Identity of the file.} The file compiled in the final build was compared with the file reported here and
  found identical.
\item \emph{Fidelity of statements.} The formal axioms and theorem statements were compared with the paper by the
  assistant; this comparison is to be repeated by the author before the code is released.
\end{itemize}

The remaining warnings concern style and the use of renamed library lemmas; they do not affect correctness. The Lean
source and the build log are available in the repository
\url{https://github.com/hirbod-assa/loss-ratio-lean}, which contains the Lean project (Lean and Mathlib
version~4.36.0-rc1) together with instructions for rebuilding it with \texttt{lake build}.
